\section{Discussion and Limitations}
\label{sec:discussion}

\subsection{Predictive Value Versus Interpretability}

The added complexity of \tabicl was not justified. Its advantage over XGBoost and the main-effects
EBM lay within sampling noise on every metric, and EBM matched or slightly exceeded it at the
operating point (recall 0.512 vs.\ 0.507; precision 1.000 vs.\ 0.991) while being transparent by
construction. In line with \citet{rudin2019stop}, the EBM would be the preferable basis for any further
development, with logistic regression as a close alternative: the EBM offers statistically
indistinguishable performance, explanations that are the model itself rather than an approximation
of it, no sampling instability and far cheaper inference (0.007\,s versus 10.8\,s to score the 545
test students on a CPU). The black-box model nonetheless served its purpose in this study, as its
opacity is precisely what made the faithfulness, stability and human evaluations necessary.

\subsection{Interpreting the Near-Perfect Precision}

The near-perfect precision does not indicate an excellent screening tool. It arises because the top
20\% is filled with students whose records already show disengagement (Section~\ref{sec:who}); such
students are easy to rank and could largely be found with a transcript query, and some may already
be beyond the reach of a meeting. For an officer, the list therefore adds an ordering rather than
new information. The more consequential group is the 105 unflagged validation dropouts, 93 of whom
passed at least one unit. A simple rule could handle the obvious cases and reserve meetings for these
less visible students, whose risk is harder to see in the transcript alone.

\subsection{Implications of the Explanation Audit}

The explanations were faithful under the deletion test and stable under minimal grade
perturbations, yet the exact list of top reasons often changed when the explanation was recomputed
with a different seed or comparison group. The bars should accordingly be read as an indication of
what broadly mattered rather than as a precise ranking. Seed instability may partly reflect the
one-permutation budget, the minimum that shap allows (a larger budget was not tested). Background
instability, by contrast, reflects a substantive choice of which students count as ``typical'' and
would not disappear with more computation. For this reason, the local research interface
(Appendix~\ref{sec:app-interface}) displays each student's stability alongside the reasons.

\textbf{Fee status and recourse.} Tuition fees ranked third by mean SHAP and second by permutation
importance, and clearing a
fee or debt record would remove the flag for 17 students. Within the model, the arrow runs from
settled fees to a lower score, but this must not be read as ``paying fees prevents dropout''. The
causal arrow may well point the other way: students who already plan to leave, or who face financial
difficulty, may stop paying, which makes unpaid fees a symptom rather than a cause. As a score change
is not a causal effect~\citep{karimi2021recourse}, an unpaid fee is better treated as a prompt for a
conversation.

\textbf{Hypothetical academic support.} First-semester progress was the dominant explanation group,
and academic support was also the help participants suggested most often (Section~\ref{sec:qual}).
Students who passed few or no units could therefore be offered tutoring, academic advising or a
structured study plan, with more targeted assistance for those who passed none. Such support may help
later grades and course completion, but its effect on dropout was not tested in this study. As
first-semester results are already fixed at the prediction point, these suggestions represent
possible human responses rather than changes that would alter the model's flag.

\subsection{Interpreting the Human Study}

The human study tested whether people could predict the model's flag, not whether explanations
improve support decisions. The task was easy without charts, as flagged profiles were visibly weak,
and for Group~B it partly became a chart-reading exercise, since adding up the bars largely
reconstructs the flag (Section~\ref{sec:audit}). With accuracy near the ceiling, the absence of a
gain cannot be taken as proof that explanations lack value. Two warnings nevertheless emerge. For
C01, an unflagged student, accuracy fell from 60\% with the profile alone to 20\% with the
explanation chart. This suggests that participants may have confused contributions increasing risk
relative to the reference group with crossing the cohort's top-20\% outreach cut-off; the study did
not directly test whether displaying the cut-off would prevent this confusion. The higher
confidence of the chart group, without any gain in accuracy, also signals a risk of over-reliance.

\subsection{Protected Attributes, Fairness and Deployment Boundaries}

Whether gender, age or nationality should influence who receives a meeting is contested. Retaining
these attributes reflects real differences in base rates, while a model blind to them can still be
unfair through proxies such as application route and course. On validation, however, removing them
cost almost nothing and narrowed the gender recall gap (Section~\ref{sec:sensitivity}), and
allocating outreach on the basis of who a student is, rather than how the student is performing,
would be difficult to justify. The position adopted here is therefore to \emph{audit} with protected
attributes but not to \emph{prioritise} with them. The subgroup results are descriptive and do not
constitute a fairness certificate.

\subsection{Threats to Validity}

The evidence comes from one Portuguese institution without external or temporal validation. The
study tests a single foundation model (TabPFN was not compared), assumes rather than measures the
20\% capacity, excludes enrolled students and rests on one test set of 545 students, whose intervals
are wide relative to the differences between models. The full explanation audit followed the final
test, using a refit that matches the frozen model, and its results depend on the SHAP budget and the
25-row background. The human study used volunteers and ten selected cases; its comparisons are
exploratory, and its qualitative coding was not independently replicated.
Finally, none of the results identify causes of dropout or effects of interventions.
